% !TEX program = lualatex
% !TeX program = lualatex
\documentclass[a4paper,11pt]{ltjsarticle}

\usepackage{geometry}
\usepackage{graphicx}
\usepackage{amsmath,amssymb}
\usepackage{hyperref}

\geometry{margin=25mm}
\hypersetup{colorlinks=true,linkcolor=blue,urlcolor=blue,citecolor=blue}

\title{タイトル}
\author{著者名}
\date{\today}

\begin{document}
\maketitle

\begin{abstract}
概要をここに書きます。目的・手法・結果を2〜3文でまとめます。
\end{abstract}

\section{はじめに}
ここに本文を書きます。段落の最初は全角スペースで始めても良いです。

\subsection{箇条書き}
\begin{itemize}
  \item ポイント1
  \item ポイント2
\end{itemize}

\subsection{数式}
式\eqref{eq:sample}のように参照できます。
\begin{align}
  E &= mc^2 \label{eq:sample}
\end{align}

\subsection{図と表}
図\ref{fig:sample}、表\ref{tab:sample}のように参照できます。
\begin{figure}[t]
  \centering
  \fbox{\rule{0pt}{120pt}\rule{0.9\linewidth}{0pt}}
  \caption{図の例（画像を入れる場合は \texttt{\\includegraphics} を使います）}
  \label{fig:sample}
\end{figure}

\begin{table}[t]
  \centering
  \caption{表の例}
  \label{tab:sample}
  \begin{tabular}{lrr}
    \hline
    項目 & 値1 & 値2 \\
    \hline
    A & 1 & 2 \\
    B & 3 & 4 \\
    \hline
  \end{tabular}
\end{table}

\section{まとめ}
結論や今後の課題を書きます。

\begin{thebibliography}{9}
  \bibitem{sample}
    著者名, \emph{文献タイトル}, 出版社, 2026.
\end{thebibliography}

\end{document}
